\section{Full Synthesis-Estimate Results}\label{app:results}

This section collects the complete synthesis-estimate evidence behind the analysis of
Sec.~\ref{app:scoring} (all resource and timing numbers are Vitis HLS \emph{csynth
estimates}, not post-place-and-route measurements). Every number is copied verbatim from the
committed canonical ablation table (\texttt{docs/ablations/canonical/TABLE.md}),
which is generated, not hand-maintained, from real Vitis HLS 2025.2 runs on part
\texttt{xc7z020clg400-1} at a 10\,ns clock---or, for the pre-fix records
preserved below, from their (locally regenerated) pre-fix run logs (e.g.\ the \texttt{stencil3\_001}
best under the old objective keeps FF~54; the canonical table's FF~43 candidate
was rejected by that objective).

The numbers below use two different objectives, and the distinction matters. The
repair demonstration (\S\ref{app:repair}) does not depend on the objective at all.
The optimize kernels re-baselined under the current default (\texttt{interval\_max}
throughput with \texttt{satisfice\_then\_area}) are in Sec.~\ref{app:scoring:rebaseline}.
The before/after tables (\S\ref{app:ppa}, \S\ref{app:matmul}) and the full suite table
(\S\ref{app:suite}) instead keep their original numbers under the earlier lexicographic
objective, and each caption flags them as historical. Where a historical per-kernel
record disagrees with the canonical evidence table, the canonical table
(\texttt{docs/ablations/canonical/TABLE.md}) is authoritative.

\subsection{Repair Metrics}\label{app:repair}

Task \texttt{vadd\_buggy\_001} is a vector add with a planted wrong operator
(\texttt{c[i] = a[i] - b[i]}, should be \texttt{+}). The agent repaired it in two
steps with a single LLM call at \$0 cost (local model). Table~\ref{tab:repair}
reports the full repair metrics.

\begin{table}[h]
\centering
\small
\begin{tabular}{@{}ll@{}}
\toprule
Metric & Value \\
\midrule
Repaired & true (winner \texttt{cand\_0001}) \\
Steps & 2 \\
LLM calls & 1 \\
Tokens (prompt / completion / total) & 2054 / 346 / 2400 \\
Budget spent (csim / csynth / prop$^{\dagger}$) & 2 / 0 / 1 \\
Cost & \$0 (local model) \\
\bottomrule
\end{tabular}
\caption{Repair metrics for \texttt{vadd\_buggy\_001}: a one-character functional
fix found in a single local-LLM call at zero cost. $^{\dagger}$\texttt{prop} counts
provider \emph{proposals} (here one real LLM call), not tokens; recipe proposals cost
0 tokens.}
\label{tab:repair}
\end{table}

\subsection{Optimize QoR Before/After (historical, pre-fix objective)}\label{app:ppa}

The three 1-D reduction tasks were driven by the deterministic \texttt{recipe}
provider and spent zero LLM tokens. The winning move in every case was a cyclic
array partition on the kernel input that unblocked pipelining to II=1, kept only
when the candidate re-passed csim and strictly improved the score.
Table~\ref{tab:ppa} reports QoR before $\rightarrow$ after under the \emph{pre-fix}
lexicographic objective (the corrected-objective re-baseline of these kernels is in
Sec.~\ref{app:scoring:rebaseline}). \texttt{mac8\_001} gets
faster \emph{and} smaller, and \texttt{unroll8\_001} is the headline latency win
($\approx$7.8$\times$, 1026$\rightarrow$132) while also shrinking LUT and FF. For
\texttt{stencil3\_001} the LUT rose (193$\rightarrow$435) and Fmax fell, but II and
latency both improved; the (then-default) lexicographic score prioritizes II then
latency over LUT, so the design is correctly accepted under that objective.

\begin{table}[h]
\centering
\small
\resizebox{\columnwidth}{!}{%
\begin{tabular}{@{}lccccc@{}}
\toprule
Task & II & Latency (worst) & LUT & FF & Fmax (MHz) \\
\midrule
\texttt{mac8\_001}     & 4$\rightarrow$1 & 1026$\rightarrow$259 & 369$\rightarrow$315 & 153$\rightarrow$126 & 144.45$\rightarrow$144.45 \\
\texttt{stencil3\_001} & 2$\rightarrow$1 & 514$\rightarrow$259  & 193$\rightarrow$435 & 95$\rightarrow$54   & 196.43$\rightarrow$149.81 \\
\texttt{unroll8\_001}  & 8$\rightarrow$1 & 1026$\rightarrow$132 & 727$\rightarrow$597 & 451$\rightarrow$368 & 144.45$\rightarrow$144.45 \\
\bottomrule
\end{tabular}}
\caption{Optimize QoR before $\rightarrow$ after using the recipe provider at 0
tokens. All three kernels reach II=1 via \texttt{ARRAY\_PARTITION cyclic factor=8
dim=1} on \texttt{in}; every accepted candidate re-passed C-simulation before
synthesis metrics were used.}
\label{tab:ppa}
\end{table}

\subsection{Generalization to a 2-D Nested-Loop Kernel (historical, pre-fix objective)}\label{app:matmul}

The 1-D fixtures above are all reductions. \texttt{matmul\_001} (8$\times$8 integer
matmul, triple nested loop) shows the same loop works on a canonical 2-D HLS kernel,
here under the \emph{original} lexicographic objective (see
Sec.~\ref{app:scoring:rebaseline} for the same kernel re-baselined under the current
\texttt{satisfice\_then\_area} default). Table~\ref{tab:matmul-prefix} reports the
before $\rightarrow$ after: the loop accepted candidate~1 (II 4$\rightarrow$1, LUT
706$\rightarrow$573, DSP 6$\rightarrow$3) and rejected the two non-improving
follow-ups. Honest nuance: \texttt{latency\_worst} \emph{rose} 260$\rightarrow$518
even as II improved, because the lexicographic score ranks II ahead of latency, so an
initiation-interval win is taken even at a single-call-latency cost. The optimizer
generalizes beyond toy reductions and never sacrifices correctness.

\begin{table}[h]
\centering
\small
\resizebox{\columnwidth}{!}{%
\begin{tabular}{@{}lcccccc@{}}
\toprule
Task & II & Latency (worst) & LUT & FF & DSP & Fmax \\
\midrule
\texttt{matmul\_001} & 4$\rightarrow$1 & 260$\rightarrow$518 & 706$\rightarrow$573 & 579$\rightarrow$612 & 6$\rightarrow$3 & 144.45$\rightarrow$144.68 \\
\bottomrule
\end{tabular}}
\caption{Generalization to a 2-D nested-loop kernel (\texttt{matmul\_001}, 3 steps,
winning pragmas \texttt{PIPELINE II=1} (inner) + \texttt{ARRAY\_PARTITION cyclic
factor=N dim=1} on \texttt{B}), under the original lexicographic objective.}
\label{tab:matmul-prefix}
\end{table}

\subsection{Token Consumption by Phase}\label{app:tokens}

The budgeted setting requires accounting token spend alongside tool calls.
Table~\ref{tab:tokens}
reports token consumption per phase across the demonstrated suite: the whole suite
(one real LLM repair plus the recipe-driven QoR optimizations) cost 2400 local
tokens (\$0) and stayed well inside the per-task budgets (csim limit 20, csynth 10,
llm\_calls 30).

\begin{table}[h]
\centering
\small
\begin{tabular}{@{}lccc@{}}
\toprule
Phase & Prompt & Completion & Total \\
\midrule
repair   & 2054 & 346 & 2400 \\
optimize & 0    & 0   & 0 \\
\textbf{all} & \textbf{2054} & \textbf{346} & \textbf{2400} \\
\bottomrule
\end{tabular}
\caption{Token consumption by phase across the four demonstrated tasks of
Table~\ref{tab:suite}: 2400 total tokens, 42 total tool calls (sum of
\texttt{budget.spent} across phases). The optimize phase spends zero tokens
because it is driven by the deterministic recipe library. \texttt{runs/SUITE.md}
additionally logs two deterministic mock-repair fixtures (0 tokens, 6 further
tool calls), excluded here because they exercise no LLM.}
\label{tab:tokens}
\end{table}

\subsection{Full Suite Table}\label{app:suite}

Table~\ref{tab:suite} aggregates every phase log into \texttt{runs/SUITE.md}
(regenerated locally; \texttt{runs/} is gitignored, and the canonical committed copy
of these numbers is under \texttt{docs/ablations/canonical/}). These rows are under
the \emph{pre-fix} lexicographic objective. The \texttt{conv2d\_001} optimize row is
omitted from this table because the locally regenerated \texttt{SUITE.md} left its
cells as placeholders; its phase log \emph{is} committed
(\texttt{docs/ablations/conv2d\_001\_optimize.json}) and records a no-improvement
outcome under the pre-fix objective---baseline kept at \texttt{interval\_max}~191,
latency~190, 871~LUT, 903~FF over 2 steps, from \emph{mock-provider} proposals at
0 tokens, so it carries no LLM evidence. (Under the corrected objective the
recipe arm does improve this kernel; see \texttt{docs/ablations/canonical/}.)
Suite-table caveat: \texttt{run\_suite.py} renders
the latency column from \texttt{latency\_best}, while the per-task analysis uses
\texttt{latency\_worst}; for these kernels the two coincide.

\begin{table*}[t]
\centering
\small
\resizebox{\textwidth}{!}{%
\begin{tabular}{@{}llllllllllll@{}}
\toprule
Task & Phase & Repaired & Improved & II (base$\rightarrow$best) & Latency (base$\rightarrow$best) & LUT (base$\rightarrow$best) & FF (base$\rightarrow$best) & Fmax & Steps & Tokens (P/C/total) & Budget (csim/csynth/prop) \\
\midrule
\texttt{mac8\_001}       & optimize & --   & True & 4$\rightarrow$1 & 1026$\rightarrow$259 & 369$\rightarrow$315 & 153$\rightarrow$126 & 144.45 & 4 & 0/0/0 & 5/5/4 \\
\texttt{stencil3\_001}   & optimize & --   & True & 2$\rightarrow$1 & 514$\rightarrow$259  & 193$\rightarrow$435 & 95$\rightarrow$54   & 149.81 & 3 & 0/0/0 & 4/4/3 \\
\texttt{unroll8\_001}    & optimize & --   & True & 8$\rightarrow$1 & 1026$\rightarrow$132 & 727$\rightarrow$597 & 451$\rightarrow$368 & 144.45 & 4 & 0/0/0 & 5/5/4 \\
\texttt{vadd\_buggy\_001} & repair  & True & --   & --              & --                   & --                  & --                  & --     & 2 & 2054/346/2400 & 2/0/1 \\
\bottomrule
\end{tabular}}
\caption{Full suite table (\texttt{runs/SUITE.md}). The \texttt{prop} column counts
provider proposals (recipe or LLM), not LLM tokens---recipe proposals cost 0 tokens
(see the Tokens column). The \texttt{conv2d\_001}
optimize row is omitted because \texttt{SUITE.md} rendered placeholder cells for it;
its committed phase log records a no-improvement outcome;
\texttt{vadd\_001} (Gate-0 toolchain proof) has no phase log; the
mock-provider repair fixtures \texttt{vadd\_offbyone\_001} and
\texttt{scale\_wrongop\_001} are logged in \texttt{SUITE.md} but omitted here
(deterministic mock patches, 0 tokens---no LLM evidence).}
\label{tab:suite}
\end{table*}
